\section{Before Starting the Campaign}

Here are some tips and suggestions intended for the GM. Following them is not mandatory, but it might help you prepare and run the game.

\subsection{Required Knowledge}

TODO; \lipsum[1][1-5]

\subsection{Read the Adventure}

You, as a GM, are expected to know the overall story of the campaign, so you could present the events with greater ease to your players, and, if a situation calls for it, modify them to your table's needs. Please familiarize yourself with most important NPCs, their goals and beliefs.

\subsection{Players’ Expectations}
\label{ssec:playersExpectations}

Inform the players that TODO; \lipsum[1][1-5]

On top of that, please discuss with your players following topics:
\begin{itemize}
	\item How serious should the adventure be?
	\item What topics players would prefer to avoid during the game?
	\item Is there anything the players would like to see in the game?
	\item Are the players comfortable with their characters dying?
	\item Are the players comfortable with characters’ romance, nudity or other sexual content?
	\item If you are of age - should your group allow alcoholic drinks?
	\item What rule variants should be in effect?
	\item How often should your group meet and how long should a single game session take?
	\item How should you handle situations when one player cannot attend a game session?
\end{itemize}

Gathering answers to those questions is not a requirement, but it will help you create a comfortable and enjoyable atmosphere for everyone.


\subsection{Create Player Characters}

Unless your players have their characters ready after previous adventures, they must create a new party. This book contains several pre-made character sheets you can use in your game (see \hyperref[chap:exampleSheets]{Appendix \ref{chap:exampleSheets}}.), but you should encourage players to create more personalized heroes for the story.

While creating new characters, you might want to follow those guidelines:
\begin{itemize}
	\item Encourage Human Ancestry, unless the player is familiar with more complex Singer rules and culture.
	\item Encourage TODO culture or at least an expertise in it, as most of the story will take place in that area.
	\item Share \hyperref[backstoryIdeas]{\textbf{Character Backstory Ideas table}} with your players for possible inspirations for characters, their motivations and goals for the story.
	\item If a player plans playing as Knight Radiant, make sure you know what spren to introduce in the story.
	\item TODO - other
\end{itemize}

\subsection{Character Relationships}

Allow players do determine how their characters know each other and what is the reason they are traveling together. Establishing characters’ relations before the start of the adventure will make it less likely the heroes will be motivated to split the party and follow their own goals.

You can add more details to character’s relationship by suggesting each player to ask about some specific details two other characters. Players can think of their own questions they want to ask each other or they can choose something from the list of below:

\begin{itemize}
	\item TODO
	\item TODO
	\item TODO
	\item TODO
\end{itemize}

\newpage
\subsection{Character Backstory}

Encourage players to write a short character backstory to better root them in the world and the story. Consider if any NPCs from those backstories could make an appearance in this adventure and how they can enrich the story. You can modify scenes described in this book to include those NPCs, for example:
\begin{itemize}
	\item TODO
	\item TODO
	\item TODO
\end{itemize}

